\chapter{Wymagania prawne}\label{ch:prawo}

\begin{uwaga}Przepisy i formularze się zmieniają. Przed każdą misją sprawdź aktualne wymagania na stronach PAŻP i w tekście rozporządzenia SERA.\end{uwaga}

\begin{sekcja}{Klasyfikacja balonu (SERA)}
Loty bezzałogowych balonów wolnych reguluje rozporządzenie wykonawcze Komisji (UE) nr 923/2012 (SERA), Dodatek~2~\cite{sera}. Klasa balonu określa wymagania. Określ klasę przed wypełnieniem zgłoszenia.

\renewcommand{\arraystretch}{1.35}
\begin{tabularx}{\linewidth}{@{}P{20mm}L@{}}
\toprule
\textbf{Klasa} & \textbf{Kryteria}\\\midrule
Lekki & Łączna masa ładunku mniejsza niż \SI{4}{kg} (jeżeli nie spełnia kryteriów balonu ciężkiego).\\
Średni & Dwa lub więcej pakietów o łącznej masie od \SI{4}{kg} do mniej niż \SI{6}{kg}.\\
Ciężki & Co najmniej jeden warunek: łączna masa \SI{6}{kg} lub więcej; jeden pakiet \SI{3}{kg} lub więcej; pakiet \SI{2}{kg} lub więcej o gęstości powierzchniowej większej niż \SI{13}{g/cm^2}; linka lub urządzenie podwieszenia, które wymaga siły \SI{230}{N} lub więcej, aby oddzielić ładunek od balonu.\\
\bottomrule
\end{tabularx}

Gęstość powierzchniowa = masa pakietu [g] / pole najmniejszej ściany pakietu [cm²].

\textbf{Zalecenia projektowe:}
\begin{itemize}
  \item Projektuj zestaw jako balon \textbf{lekki}. Wymagania dla balonu ciężkiego (dwa niezależne systemy zakończenia lotu, transponder, oświetlenie, raportowanie pozycji) są trudne do spełnienia dla zespołu studenckiego.
  \item Użyj linki, która zrywa się przy sile mniejszej niż \SI{230}{N} -- w praktyce linka o średnicy około \SI{1}{mm}. Sprawdź wytrzymałość w karcie katalogowej albo wykonaj odpowiedni test (np. zerwij próbkę linki, mierząc siłę wagą bagażową).
  \item Podziel ładunek na pakiety o masie poniżej \SI{2}{kg}.
\end{itemize}
\end{sekcja}

\begin{sekcja}{Zgłoszenie lotu do PAŻP}
Polska Agencja Żeglugi Powietrznej (PAŻP) przyjmuje zgłoszenia sondowań atmosfery i publikuje ostrzeżenie dla pilotów (NOTAM)~\cite{pazp}.

\textbf{Formularz:} formularz zgłoszenia sondowania atmosfery (ASM1):\\
\url{https://www.pansa.pl/wp-content/uploads/2025/07/Formularz_sonda_v1.pdf}\\
Lista aktualnych formularzy: \url{https://www.pansa.pl/asm1/formularze-zgloszeniowe/}

\textbf{Procedura:}
\begin{enumerate}
  \item Sprawdź strefy przestrzeni powietrznej w miejscu startu i na przewidywanej trasie lotu (rozdział~\ref{ch:planowanie}, sekcja~\ref{sec:strefy}).
  \item Pobierz aktualny formularz (adres wyżej). Sprawdź, czy PAŻP nie opublikowała nowej wersji.
  \item Wypełnij \textbf{dwa formularze}: dla miejsca zasadniczego i zapasowego. W każdym podaj termin zasadniczy i zapasowy (rozdział~\ref{ch:planowanie}).
  \item Wyślij formularze e-mailem do działu ASM1: \texttt{asm1@pansa.pl}, co najmniej 14 dni przed lotem. Minimalny termin podaje strona PAŻP -- sprawdź go przed każdą misją. Dział ASM1 pracuje w dni robocze.
  \item Zapisz w karcie misji (F0) numer zgłoszenia i numer telefonu do służby ruchu lotniczego z odpowiedzi PAŻP.
  \item Przed startem sprawdź w AIS Polska, czy NOTAM dla misji jest opublikowany.
  \item W dniu startu miej numer do służby ruchu lotniczego pod ręką. Gdy zmienisz plan (godzina, odwołanie), zadzwoń natychmiast.
\end{enumerate}
\end{sekcja}

\begin{sekcja}{Dane w formularzu PAŻP}
\renewcommand{\arraystretch}{1.3}
\begin{tabularx}{\linewidth}{@{}P{26mm}L@{}}
\toprule
\textbf{Sekcja} & \textbf{Dane}\\\midrule
Organizator & Imię i nazwisko, telefon, e-mail.\\
Termin & Termin zasadniczy i zapasowy: data, godziny (UTC), planowana godzina startu.\\
Lokalizacja & Nazwa miejscowości, współrzędne WGS-84 (stopnie, minuty, sekundy).\\
Balon & Kolor, ilość, klasa (lekki/średni/ciężki), masa powłoki, wysokość sondowania, prędkość wznoszenia i opadania, średnica przy starcie i na wysokości pęknięcia.\\
Ładunek & Kolor i średnica spadochronu, kolor ładunku, częstotliwość nadajnika, masa ładunku ze spadochronem, gęstość powierzchniowa (gdy pakiet ma \SI{2}{kg} lub więcej).\\
Opis & Krótki opis przedsięwzięcia.\\
\bottomrule
\end{tabularx}

Dane o średnicy, wysokości i prędkości weź z kalkulatora (rozdział~\ref{ch:planowanie}).
\end{sekcja}

\begin{sekcja}{Licencja radiowa}
Nadajniki trackera (APRS na \SI{144.8}{MHz}, RTTY na \SI{434}{MHz}) wymagają pozwolenia radiowego dla służby amatorskiej. Operator z licencją krótkofalowca odpowiada za nadajnik i udostępnia swój znak wywoławczy.

\begin{itemize}
  \item Jak zdobyć licencję krótkofalowca w Polsce (egzamin, pozwolenie, znak): przewodnik Hackerspace Kraków~\cite{faqradiowe}: \url{https://github.com/HackerspaceKRK/faq-radiowe}.
  \item Sprawdź, czy warunki pozwolenia dopuszczają pracę z balonu. Niektóre kraje zabraniają APRS w powietrzu.
  \item Moduły ISM (np. \SI{434}{MHz}, LoRa) mają limity mocy i cyklu pracy. Sprawdź je dla użytego pasma.
\end{itemize}

\begin{wskazowka}Zdobądź licencję w zespole z wyprzedzeniem -- egzamin i wydanie pozwolenia trwają tygodnie.\end{wskazowka}
\end{sekcja}

\begin{sekcja}{Znak wywoławczy misji}
Każda misja nadaje pod znakiem wywoławczym konkretnej osoby albo klubu. Właściciel znaku odpowiada za nadajnik.

Kolejność wyboru znaku:
\begin{enumerate}
  \item \textbf{Znak klubowy SimLE} -- docelowo koło ma własny klubowy znak krótkofalarski.\par
  {\raggedright\pole[45mm]{Znak klubowy:}\par\pole[45mm]{Osoba odpowiedzialna:}\par}
  \item \textbf{Znak członka zespołu} z licencją krótkofalowca.
  \item \textbf{Znak udostępniony przez osobę spoza zespołu} -- tylko za jej wyraźną zgodą.
\end{enumerate}

\begin{uwaga}[Znak SR0FLY]Znak SR0FLY, używany w misji z 27.02.2022, należy do Macieja Jakimca (Fundacja Copernicus Project, nearspace.pl). Przed każdym użyciem tego znaku \textbf{bezwzględnie} skontaktuj się z nim i poproś o zgodę:\\
\url{https://nearspace.pl/pl/cms/o-nas-30.html}\\
\url{https://nearspace.pl/pl/contact.html}\end{uwaga}

Zapisz w karcie misji (F0): znak, SSID, właściciela znaku i datę jego zgody.
\end{sekcja}

\begin{sekcja}{Ubezpieczenie OC}
Wykup ubezpieczenie odpowiedzialności cywilnej na lot.
\begin{itemize}
  \item Zespół KSat e.V. z Uniwersytetu w Stuttgarcie, który prowadzi własne misje balonowe (projekt BUBBLE)~\cite{ksat}, rekomenduje lotnicze ubezpieczenie OC firmy Allianz według ogólnych warunków AMU~300/06~\cite{amu}.
  \item Sprawdź, czy polisa obejmuje bezzałogowe balony wolne. Ogólne polisy lotnicze często nie wymieniają balonów wprost.
  \item W Polsce: zapytaj ubezpieczyciela uczelni albo brokera o OC organizacji z klauzulą lotniczą. Podaj klasę balonu, masę ładunku i obszar lotu.
\end{itemize}
\end{sekcja}

\begin{sekcja}{Oznakowanie ładunku}
Opisz każdy pakiet. Wystarczy pisać markerem permanentnym bezpośrednio na styropianie.
\begin{itemize}
  \item \textbf{Numer telefonu -- obowiązkowo.} Osoba, która znajdzie gondolę, musi móc zadzwonić.
  \item Informacja dla osób postronnych (zalecana): „Projekt badawczy -- nie otwierać, prosimy o kontakt”.
  \item Logo projektu (naklejka) -- zwiększa zaufanie znalazcy.
  \item Ostrzeżenia, np. o bateriach litowych.
\end{itemize}

\dwafoto{2023_gondola_bok_termometr}{Oznakowanie markerem: „Projekt badawczy”. Z boku wyprowadzony czujnik temperatury.}
{2023_gondole_po_locie}{Numer telefonu („Znalazłeś? Zadzwoń!” -- numer zamazany w tej publikacji), napis „Projekt badawczy” i naklejki z logo misji.}
\end{sekcja}
